% !PRINTED 261
\weberpage{261}

III.\ \gesperrt{Man kann, wenn $R$ von Null verschieden ist, die Indices $1, 2 \ldots m$ so anordnen, dass in der Kette der Haupt-Unterdeterminanten}
\weq{3}{R_m,\ R_{m-1} \ldots R_1,\ R_0}
\gesperrt{nicht zwei auf einander folgende Glieder verschwinden.}

Die Determinantenrelation (2) ergiebt, wenn man
\[
A = R_{k+1}, \quad (k = 1, 2 \ldots m-1)
\]
setzt, eine Gleichung zwischen drei auf einander folgenden Gliedern der Kette (3), die wir so schreiben können:
\weq{4}{R_k S_k - T_k^2 = R_{k-1} R_{k+1},}
worin $S_k$ und $T_k$ gewisse Unterdeterminanten von $R$, also ganze rationale Functionen der Co\"efficienten $a_{i,k}$ sind.

Näher bezeichnet, sind $R_k$, $S_k$, $T_k$ erste Unterdeterminanten von $R_{k+1}$, und zwar ist
\[
R_k = \frac{\partial R_{k+1}}{\partial a_{k+1,k+1}}, \quad S_k = \frac{\partial R_{k+1}}{\partial a_{k,k}}, \quad T_k = \frac{\partial R_{k+1}}{\partial a_{k,k+1}}.
\]

Eine dem Satz (3) entsprechende Anordnung der Indices wollen wir für die Folge als gewählt voraussetzen. Dann ist, wenn $R_k$ verschwindet, $R_{k-1}$ und $R_{k+1}$ von Null verschieden und (4) zeigt, dass sie entgegengesetzte Vorzeichen haben, also:

IV.\ \gesperrt{Wenn ein inneres Glied einer Kette von Haupt-Unterdeterminanten verschwindet, so haben die beiden angrenzenden Glieder entgegengesetzte Vorzeichen.}

\begin{center}
\S.\,83.\\[4pt]
\gesperrt{Anzahl der positiven und negativen Quadrate.}
\end{center}

Um die Anzahl der positiven und negativen Quadrate einer Form mit nicht verschwindender Determinante zu bestimmen, nehmen wir zunächst an, dass in der Kette der Haupt-Unterdeterminanten
\[
R_m,\ R_{m-1},\ R_{m-2} \ldots R_1,\ R_0 = 1
\]
kein Glied verschwinde. Wir können dann den Co\"efficienten $\lambda$ so bestimmen, dass die Determinante der Form
%CONT>>
